\chapter*{РЕФЕРАТ}
\addcontentsline{toc}{chapter}{РЕФЕРАТ}
\vspace{-1cm}

Дипломная работа: 65 страниц, 14 рисунков, 2 таблицы, 30 источников, 2 приложения.

\textbf{Ключевые слова:} КВАНТОВАЯ ТОЧКА, МЕТОД СИЛЬНОЙ СВЯЗИ, СУПЕРЭЛЛИПС,
НИЗКОЭНЕРГЕТИЧЕСКИЙ СПЕКТР, СУРРОГАТНАЯ АППРОКСИМАЦИЯ,
ГРЕБНЕВАЯ РЕГРЕССИЯ, МНОГОСЛОЙНЫЙ ПЕРЦЕПТРОН,
СТРУКТУРИРОВАННАЯ ПРОВЕРКА.\par

\textbf{Объект исследования} --- квантовые точки суперэллиптической формы на квадратной решётке метода сильной связи. \textbf{Предмет исследования} --- энергетический спектр, соответствующие собственные состояния дискретного гамильтониана и зависимость \(E_0\) и \(dE_1\) от размера и формы границы. \textbf{Цель работы} --- моделирование энергетического спектра таких точек и проверка физически мотивированной суррогатной аппроксимации; волновые функции рассматриваются как собственные состояния спектральной задачи.

\textbf{Методы исследования:} расчёт спектра и собственных состояний в Kwant, прямоугольная контрольная задача, проверки \(E_{\mathrm{kin}}=E_0+4\), \(1/a^2\) и Бесселя, Ridge-регрессия, малый многослойный перцептрон, схемы LOAO/LOARO и анализ остатков.

\textbf{Основные результаты.} Рассчитаны 140 геометрий: \(n=\{1.2,2.0,3.0,4.0\}\), \(a=\{24,27,30,33,36\}\), \(r_{\mathrm{AR}}\in[0.67,1.0]\). Ошибка бесселевой оценки равна примерно \(2.03\%\), максимальное отклонение проверки масштабирования при фиксированной форме --- \(2.12\%\). МП с физически мотивированными признаками успешен в 10 из 16 ячеек: 3 из 8 для LOAO и 7 из 8 для LOARO. Заданный критерий не выполнен, поэтому основной моделью остаётся Ridge. Относительная MAE Ridge составляет 2.01--2.60\% от \(E_{\mathrm{kin}}\) для \(E_0\) и 0.78--1.15\% от \(dE_1\). Глобальное объяснение остатков простыми диагностиками границы не подтверждено; отдельный сигнал найден только для \(n=1.2\) (\(\rho\approx -0.416\), \(p\approx0.013\)). \textbf{Ограничения.} Результаты относятся только к исследованной сетке параметров и фиксированным классам \(n\); суррогатная модель не заменяет прямой расчёт Kwant.
\par
